\scanpage{41}
% sheet: reused_typed
% hand: typed
% L1 | top centre, in reading orientation
\mspage{-3.13-}

% L2–3 | upper
The key property of stable geometric points turns out to be the following:

% L4–9 | upper and middle
\uline{Proposition} \uline{3.6}: Let $x$ be a geometric point of $(\handwritten{\mathbb P}_n)^{m+1}$. Then $x$ is stable if andonly if there are integers $N\geq N_0>0$ and there are monomials $P_0,\ldots,P_m$ in the expressions $D_{\handwritten{\alpha_0},\ldots,\handwritten{\alpha_n}}$ such that, for $0\leq i\leq m$:

(i) $P_i$ is a section of $L_0^N\handwritten{\otimes}\cdots\handwritten{\otimes}L_i^{N-N_0}\handwritten{\otimes}\cdots\handwritten{\otimes}L_m^N$

(ii) $P_i(x)\ne0$

% L10–14 | middle
\uline{Proof}: Let $E=(\handwritten{\mathbb R})^{m+1}$ be Euclidean space of dimension $m+1$. To each $D_{\handwritten{\alpha_0},\ldots,\handwritten{\alpha_n}}$ associate a point $d_{\handwritten{\alpha_0},\ldots\handwritten{\alpha_n}}\in E$ as follows:
\[
\begin{aligned}
d_{\handwritten{\alpha_0},\ldots\handwritten{\alpha_n}}&=(\unsure{\epsilon|\varepsilon}_0,\epsilon_1,\ldots,\epsilon_m)\\
\epsilon_i&=0\quad\text{if }\ i\handwritten{\notin\{}\handwritten{\alpha_0},\ldots\handwritten{\alpha_n}\handwritten{\}}\\
\epsilon_i&=1\quad\text{if }\ i\handwritten{\in\{}\handwritten{\alpha_0},\ldots\handwritten{\alpha_n}\handwritten{\}}
\end{aligned}
\]

% L15–20 | lower
Then $d$ determines which invertible sheaf $D$ is a section of. More generally, to any monomial $P$ is the $D$'s, associate the corresponding sum of the $d$'s. Now define $C_x$ to be the convex cone spanned by those points $d_{\handwritten{\alpha_0},\ldots,\handwritten{\alpha_n}}$ such that $D_{\handwritten{\alpha_0},\ldots,\handwritten{\alpha_n}}(x)\ne0$, i.e. such that the components $x^{(\handwritten{\alpha_0})},\ldots,x^{(\handwritten{\alpha_n})}$ do not lie in a hyperplane. Then it is easy to see that:
% archive sheet number: 40

